\chapter{A running example: an order with no first element}
\label{tx:ch:running-order}
We now apply the consequence-space viewpoint to one concrete family.
The equations describe a finite total order in which every element has
a selected predecessor. A finite order has a first element, so each
instance is inconsistent. The question is how much polynomial degree
is needed to expose that contradiction.

For the running example fix a finite field $K$ of characteristic $p$.
The chain and restriction arguments below work over every field; the
final linear bound also holds over any fixed field of positive characteristic.

The proof has three stages. First, a restriction to an elementary
induction chain shows that NS degree tends to infinity. Second, the
uniform NS--PC comparison proved in Chapter~\ref{tx:sec:comparison}
shows that PC degree also tends to infinity. Third, the degree-preserving
lifting theorem, constructed in the intervening chapters and assembled
in Chapter~\ref{tx:chap:assemble-transfer}, upgrades this qualitative
information to a linear lower bound. The chain obstruction and the
restriction are proved here in full. The two general results are stated
at their point of use, so the reader can see exactly what their later
proofs must deliver.

\section{The equations and their finite proofs}
For this chapter write $[n]=\{1,\ldots,n\}$; this is just a
renaming of the book's universe. There are two sorts of variables:
\[
 L_{ij}\quad\text{means that }i\text{ precedes }j,\qquad
 F_{ij}\quad\text{means that the selected predecessor of }i\text{ is }j.
\]
Let $\mathcal O_n$ consist of the following polynomials, all interpreted
as equations equal to zero:
\begin{subequations}\label{tx:eq:running-order-axioms}
\begin{align}
 L_{ij}^2-L_{ij},\quad F_{ij}^2-F_{ij}
       &\qquad(i,j\in[n]),\label{tx:eq:order-Boolean}\\
 L_{ii}&\qquad(i\in[n]),\label{tx:eq:order-irreflexive}\\
 L_{ij}L_{ji},\quad(1-L_{ij})(1-L_{ji})
       &\qquad(i\ne j),\label{tx:eq:order-total}\\
 L_{ij}L_{jk}(1-L_{ik})
       &\qquad(i,j,k\in[n]),\label{tx:eq:order-transitive}\\
 \sum_{j=1}^nF_{ij}-1
       &\qquad(i\in[n]),\label{tx:eq:order-function-total}\\
 F_{ij}F_{ik}&\qquad(j\ne k),\label{tx:eq:order-function-collision}\\
 F_{ij}(1-L_{ji})&\qquad(i,j\in[n]).\label{tx:eq:order-predecessor}
\end{align}
\end{subequations}
Here a selected predecessor means any strictly earlier element; it
need not be the immediately preceding element. Different elements may
choose the same predecessor.

The first four lines describe a strict total order. The next two select
exactly one predecessor in each row. The last requires that this
predecessor really come earlier. This is the graph-and-sum encoding of
the first-order assertion that $L$ is a strict total order and
$L(F(i),i)$ holds for every $i$.

The functionality collisions matter in positive characteristic.
Boolean values and a sum equal to one alone could permit $p+1$ selected
entries. The collisions allow at most one, and the sum then forces one.

\begin{proposition}[Finite inconsistency and finite proof degree]
\label{tx:prop:order-finite}
For every $n\ge1$ and every field $K$, the system $\mathcal O_n$ has
an NS refutation and an ordinary-PC refutation. Their minimum degrees
are finite and positive.
\end{proposition}
\begin{proof}
At a root, all variables are zero or one. The order equations give a
strict total order on the labels. Starting from any label, move to an
earlier one while possible. Transitivity and irreflexivity forbid a
repeated label, so finiteness forces this process to stop at a first
element $i$.
The function equations select a label $j$ with $F_{ij}=1$.
Equation~\eqref{tx:eq:order-predecessor} then gives $L_{ji}=1$,
contradicting the choice of $i$. The same argument applies over every
extension field.

Boolean refutation completeness
(Proposition~\ref{tx:prop:boolean-completeness}) constructs a finite
NS certificate, and its simulation gives a finite PC proof. No input
is a nonzero constant, so the degree-zero consequence space is zero.
Neither refutation can have degree zero.
\end{proof}

At $n=1$ the contradiction is already visible in six ordinary-PC lines:
\begin{center}
\begin{tabular}{@{}cllc@{}}
\toprule
Line & Polynomial & Reason & Degree\\
\midrule
1 & $L_{11}$ & input & 1\\
2 & $F_{11}-1$ & input & 1\\
3 & $F_{11}(1-L_{11})$ & input & 2\\
4 & $F_{11}L_{11}$ & multiply line 1 by $F_{11}$ & 2\\
5 & $F_{11}$ & add lines 3 and 4 & 1\\
6 & $1$ & subtract line 2 from line 5 & 0\\
\bottomrule
\end{tabular}
\end{center}
For larger $n$ there is no designated first label. The order variables
must determine it, while every row insists on a predecessor.

\subsection{Why this is one uniform family}
Relabelling sends $L_{ij}$ and $F_{ij}$ to the variables with the
relabelled ordered pair. Every input has an index using at most three
labels, and all coefficients follow one fixed rule. For example, the
row with index $i$ has coefficient one on every $F_{ij}$ and constant
coefficient minus one, regardless of $n$.

Safe template parameters are
\begin{equation}\label{tx:eq:order-parameters}
                  k=2,\qquad s=3,\qquad D=3,\qquad n_0=9.
\end{equation}
There are $2n^2$ variables, but one variable uses at most two labels.
The target $1$ is fully invariant and has degree zero. At proof degree
$d$ the invariant-target support budget is therefore
\begin{equation}\label{tx:eq:order-support-budget}
                              r=kd+s=2d+3.
\end{equation}
This budget bounds the bases and maps of the consequence calculation;
it does not require each proof line to mention only $2d+3$ labels.

\section{A smaller chain problem with two different proof degrees}
Here the input comes with a \emph{fixed successor chain}: the numerical
chain $1,2,\ldots,n$, its successor relation and its two endpoints
are fixed in advance. They are part of the input rule, rather than
relations to be determined by a satisfying assignment. For variables
$x_1,\ldots,x_n$, define the induction system
\begin{equation}\label{tx:eq:order-induction}
 \mathrm{IND}_n=\{1-x_1,\ x_i(1-x_{i+1})\ (1\le i<n),\ x_n,
                                      \ x_i^2-x_i\ (1\le i\le n)\}.
\end{equation}
Its meaning is sequential: the first variable is one, a one forces the
next variable to be one, and the last variable is zero.

There is an ordinary-PC refutation of degree at most two. Starting
from a derived line $P=1-x_i$, first multiply it by the single variable
$x_{i+1}$. Then form
\[
       P-x_{i+1}P+x_i(1-x_{i+1})=1-x_{i+1}
\]
by two linear combinations. Every intermediate line has degree at
most two. Iterate and add the final input $x_n$ to obtain $1$.
For $n=1$, the two endpoint inputs already sum to one.

NS behaves differently: its minimum degree is $\Theta(\log n)$.
We prove both a logarithmic lower bound and a logarithmic upper bound
in the next section. Thus the fixed chain gives a constant versus
logarithmic degree separation. The linear NS and PC bounds belong
to the original uniform order family $\mathcal O_n$, whose order is
an unknown relation. The distinction is possible because this chain family is
\emph{not} a uniform template of the kind defined in Chapter~\ref{tx:ch:uniformity}.
It singles out numerical adjacency and the last label. An arbitrary
label permutation does not preserve its input rule. Thus the uniform
NS--PC comparison cannot be applied to $\mathrm{IND}_n$ itself.

The induction system is a useful restriction of the uniform order
family. We next prove its obstruction using a finite graph of
monomials. The classical NS degree behaviour of induction is discussed
by Buss and Pitassi~\cite{BussPitassiInduction}; the needed lower bound
is proved below rather than used as an external theorem.

\section{A finite graph that witnesses the NS obstruction}
Fix an integer degree budget $b\ge0$. A multilinear monomial
$x^S=\prod_{i\in S}x_i$ is indexed by a subset $S\subseteq[n]$.
Make a graph whose vertices are the sets with $|S|\le b$. Its edges
allow the following changes:
\begin{enumerate}
\item Insert or delete label $1$.
\item Insert or delete label $i>1$ when label $i-1$ is present.
\end{enumerate}
The number of occupied labels must remain at most $b$. Call an occupied
label a pebble. The moves are reversible: the predecessor is still
present after the move, so the reverse move is legal.

\subsection{The clearance invariant}
We need a bound on how far a sequence starting with no pebbles can go.
The following auxiliary test supplies it. Write a configuration as
$a_1<\cdots<a_k$, put $a_0=0$, and call pebble $a_j$ \emph{eligible}
for removal when
\begin{equation}\label{tx:eq:order-clearance-test}
                        a_j-a_{j-1}\le2^{b-k}.
\end{equation}
A configuration is \emph{clearable} if successive eligible removals
can empty it, recalculating the threshold after each removal.
These auxiliary removals need not be legal graph moves. They give an
invariant for the actual graph moves, whose predecessor gaps are one.
In this clearance argument, $k$ counts the currently occupied pebbles;
it is independent of the variable-support parameter $k=2$ for
$\mathcal O_n$.

\begin{lemma}[Any eligible removal can come first]
\label{tx:lem:order-clearance-exchange}
Removing any eligible pebble from a clearable configuration leaves a
clearable configuration.
\end{lemma}
\begin{proof}
Induct on the number $k$ of pebbles. The assertion for $k\le1$ is
immediate. Put $h=2^{b-k}$, choose an eligible pebble $a$, and let $c$
be the first removal in one clearing order. If $a=c$, use the remainder
of that order.

Otherwise both are initially eligible. Removing either doubles the
threshold to $2h$. The other's gap stays the same unless the two
pebbles are adjacent in the configuration; in that case its gap can
become the sum of two gaps at most $h$, hence is at most $2h$. Thus either order
can remove both. The configuration with $c$ removed is clearable;
by induction it can remove $a$ next and then clear the remaining
configuration. Starting instead by removing $a$, remove $c$ next and
use that same remainder. This gives a clearing order after removing $a$.
\end{proof}

\begin{lemma}[The reachable frontier]
\label{tx:lem:order-pebble-frontier}
Every configuration reachable from the empty set using at most $b$
pebbles is clearable. In every nonempty clearable configuration of
$k$ pebbles,
\begin{equation}\label{tx:eq:order-pebble-frontier}
                     a_k\le2^b-2^{b-k}\le2^b-1.
\end{equation}
\end{lemma}
\begin{proof}
The empty configuration is clearable. For a legal deletion, the
deleted pebble has gap one, so is eligible; the exchange lemma preserves
clearability. For a legal insertion, the new pebble has gap one and
can be removed immediately by the auxiliary test. This returns to the
previous clearable configuration, proving that the new one is also
clearable. This proves the invariant along every graph path.

For the bound, induct on $k$, assigning maximum zero to the empty
configuration. A clearing order first removes an eligible pebble with
gap at most $h=2^{b-k}$. The remaining configuration is clearable and
has $k-1$ pebbles, so its maximum is at most $2^b-2h$. If the deleted
pebble was not the largest, this bounds the original maximum as well.
If it was the largest, restoring it adds at most its gap $h$ to that
maximum. Either way, $a_k\le2^b-h$. Since $k\le b$, one has $h\ge1$,
giving the final inequality.
\end{proof}

For example, two pebbles can reach label three:
\[
 \varnothing,\ \{1\},\ \{1,2\},\ \{2\},\ \{2,3\}.
\]
They cannot reach label four. The clearance test explains the limitation
without bounding the number of moves: with two occupied labels the
threshold is one, and after an eligible removal it becomes two.

\subsection{From graph components to a linear functional}
A \emph{linear functional} assigns a scalar to each coefficient vector
and respects addition and scalar multiplication. Specifying its value
on every monomial specifies it on every polynomial by linearity.
To exclude an NS identity for $1$, it suffices to find such a functional
which is zero on every permitted multiplier--input product but nonzero
on $1$. Applying it to the identity would give a contradiction.
This is a dual certificate for nonmembership in a span.

The graph supplies its weights. A difference of two adjacent monomial
coordinates receives value zero if the weights are constant along
edges, hence on connected components. We give the empty component
weight one and the other components weight zero. The proof below
checks that this assignment also handles powers and endpoint inputs.

\begin{proposition}[An elementary logarithmic NS obstruction]
\label{tx:prop:order-induction-lower}
If $n\ge2^b$, the system $\mathrm{IND}_n$ has no total-product
degree-$b$ NS refutation, over any field.
\end{proposition}
\begin{proof}
Let $\mathcal E$ be the component of the empty vertex in the graph.
By Lemma~\ref{tx:lem:order-pebble-frontier}, no vertex of $\mathcal E$
contains $n$. For a monomial $\mu$ of degree at most $b$, let $S(\mu)$
be the labels with positive exponent and define
\[
 \Lambda(\mu)=
 \begin{cases}1,&S(\mu)\in\mathcal E,\\0,&S(\mu)\notin\mathcal E.
 \end{cases}
\]
Extend $\Lambda$ linearly to $P_{\le b}$. In particular $\Lambda(1)=1$.
We check that it vanishes on every permitted monomial--input product.

The two monomials of $\mu(1-x_1)$ reduce to the same subset or to
vertices joined by a label-one edge. Their $\Lambda$ values are equal,
so their difference has value zero. For
$\mu x_i(1-x_{i+1})$, the subsets are
$S(\mu)\cup\{i\}$ and $S(\mu)\cup\{i,i+1\}$. They are identical
or adjacent, because $i$ is present. They both have size at most $b$
by the product-degree restriction, so this difference also has value
zero. The product $\mu x_n$ corresponds to a vertex containing $n$,
hence has value zero. Finally $\mu(x_i^2-x_i)$ has the same set of
positive exponents in both monomials, so again has value zero.

Every allowed NS identity is a linear combination of these products.
Applying $\Lambda$ to a supposed identity for $1$ would give $0=1$.
This contradiction proves the proposition.
\end{proof}

The functional used for the lower bound is a test on coefficient vectors, not a truth
assignment. Replacing a positive exponent by one in its definition
does not grant a free Boolean-reduction rule to a PC proof.
Indeed, the functional need not respect multiplication of derived
low-degree lines; the degree-two PC refutation above demonstrates
precisely why an NS witness need not be an ordinary-PC witness.

For a small calculation, take $b=2$ and $n\ge4$. The empty component
contains $\{2\}$ and $\{2,3\}$ but not $\{3\}$: the singleton's gap
three fails its clearance threshold two. Consequently
\[
 \Lambda(1-x_2)=0,\qquad
 \Lambda\bigl(x_3(1-x_2)\bigr)=0-1=-1\ne0.
\]
The line $1-x_2$ has a degree-two PC derivation. Its permitted
multiplication by $x_3$ is exactly an operation this NS functional does
not preserve.

\subsection{Why the chain also has logarithmic NS certificates}
The chain obstruction cannot be strengthened to a linear bound.
There is a short divide-and-conquer construction that makes this
clear without appealing to an external upper-bound theorem.

\begin{proposition}[A logarithmic NS upper bound for the fixed chain]
\label{tx:prop:order-induction-upper}
For $n\ge2$, the system $\mathrm{IND}_n$ has a total-product NS
refutation of degree at most $2+\lceil\log_2(n-1)\rceil$, over every field.
\end{proposition}
\begin{proof}
For $a<b$ write $T_{ab}=x_a(1-x_b)$. When $b=a+1$, this is
an input. For any $a<c<b$, the polynomial identity
\[
 T_{ab}=(1-x_b)T_{ac}+x_aT_{cb}
\]
combines two shorter intervals. Choose $c$ so that their lengths
differ by at most one, and recursively expand both terms until only
the original adjacent-transition inputs remain. Each recursive level
multiplies their certificates by a polynomial of degree one. An
interval of length $b-a$ needs at most
$\lceil\log_2(b-a)\rceil$ levels, so every resulting multiplier--input
product has degree at most $2+\lceil\log_2(b-a)\rceil$.
This is an NS identity in the original inputs, obtained by expanding
all intermediate transitions.

Finally use the endpoint identity
\[
                1=(1-x_1)+T_{1n}+x_1x_n.
\]
The first term is an input, the last is a degree-two multiple of the
input $x_n$, and the middle has the certificate just constructed.
The displayed degree bound follows. For $n=1$, the endpoint inputs
already sum to one with total-product degree one.
\end{proof}

Together with Proposition~\ref{tx:prop:order-induction-lower}, this
proves the claimed $\Theta(\log n)$ NS degree. The fixed successor chain
allows PC to carry a single degree-one conclusion down the chain;
an NS certificate records the original inputs, and balanced
composition keeps its growing degree logarithmic.

\section{Restricting the uniform order problem to the chain}
In this argument, a \emph{substitution} assigns a polynomial to each
variable and extends by addition and multiplication to the whole
polynomial ring. It is a ring homomorphism:
$\sigma(af+bg)=a\sigma(f)+b\sigma(g)$ and
$\sigma(fg)=\sigma(f)\sigma(g)$, where $a,b\in K$.
Thus it sends an identity $1=\sum a_ff$ to
$1=\sum\sigma(a_f)\sigma(f)$.
To obtain a certificate from the chain inputs, we must additionally
express each substituted input $\sigma(f)$ using those inputs.
The degree accounting below controls both steps. Substitution here
is a fixed-size reduction between two systems, whereas the later
lifting operator is a linear map between universe sizes.

We now relate $\mathcal O_n$ to $\mathrm{IND}_n$, for $n\ge2$.
The substitution need not respect relabelling: it will be used only
to transform an NS identity at one fixed size. Set
\begin{equation}\label{tx:eq:order-cycle-restriction}
 \sigma(F_{ij})=
 \begin{cases}
 1,&i>1\text{ and }j=i-1,\text{ or }i=1\text{ and }j=n,\\
 0,&\text{otherwise},
 \end{cases}
\end{equation}
and put
\begin{equation}\label{tx:eq:order-relation-restriction}
 \sigma(L_{ij})=
 \begin{cases}
 0,&i=j,\\
 1-x_i(1-x_j),&i<j,\\
 (1-x_i)x_j,&i>j.
 \end{cases}
\end{equation}
Every substituted variable has degree at most two. On a Boolean
assignment, the relation in \eqref{tx:eq:order-relation-restriction}
orders labels first by their Boolean value, with zero before one,
and then by their numerical label. It is therefore always a strict
total order.

The selected predecessors run around one fixed cycle. For $i>1$,
the nonzero guard becomes exactly
\[
            1-\sigma(L_{i-1,i})=x_{i-1}(1-x_i).
\]
The wraparound guard is
\begin{equation}\label{tx:eq:order-wrap-guard}
       1-\sigma(L_{n,1})=(1-x_1)+x_1x_n.
\end{equation}
This has a degree-two NS identity from the two endpoint inputs of
$\mathrm{IND}_n$. All unselected guards become zero, and the function
Boolean, totality and collision equations become zero identically.

\begin{lemma}[Boolean reduction retains a product-degree bound]
\label{tx:lem:order-Boolean-reduction}
If a polynomial $P$ of degree at most $m$ vanishes on every Boolean
assignment, it is a sum of multiples of $x_i^2-x_i$ with each nonzero
product of degree at most $m$.
\end{lemma}
\begin{proof}
Reduce powers one monomial at a time using
\[
 x_i^a-x_i=(x_i^2-x_i)(1+x_i+\cdots+x_i^{a-2})\qquad(a\ge2).
\]
Multiplying this identity by the other factors of that monomial
does not exceed its original degree. Continue until every monomial
is multilinear, recording the Boolean-axiom products. The final
multilinear remainder still vanishes on the Boolean cube, so is zero
by the induction proved in Chapter~\ref{tx:ch:consequences}.
\end{proof}

Apply this lemma to the substituted order equations, which vanish on
every Boolean assignment. A Boolean or two-label order axiom has
degree two before substitution and degree at most four afterwards.
A transitivity axiom has degree three and becomes a polynomial of
degree at most six. Irreflexivity becomes zero. Thus each source
axiom $f$ of actual degree $e$ has a substituted NS identity from
$\mathrm{IND}_n$ with product degree at most $2e$; the guards have
the stronger bound two, as just computed.

\begin{lemma}[Degree accounting for the restriction]
\label{tx:lem:order-NS-restriction}
A degree-$d$ NS refutation of $\mathcal O_n$ gives a degree-at-most-$2d$
NS refutation of $\mathrm{IND}_n$.
\end{lemma}
\begin{proof}
Write the original identity as $1=\sum_f a_ff$, omitting zero
products. If $e=\deg f$, then $\deg a_f\le d-e$. Hence
$\deg\sigma(a_f)\le2(d-e)$. Replace $\sigma(f)$ by the identity
just established, whose individual products have degree at most $2e$.
After multiplication by $\sigma(a_f)$, every new product has degree
at most $2(d-e)+2e=2d$. Summing gives the required NS refutation.
\end{proof}

\begin{corollary}[The qualitative NS seed]
\label{tx:cor:order-NS-seed}
For every $n\ge2$ and every field,
\begin{equation}\label{tx:eq:order-log-seed}
                        d_{\NS}(\mathcal O_n)>
                                      \tfrac12\log_2 n.
\end{equation}
In particular, the NS minimum degrees tend to infinity.
\end{corollary}
\begin{proof}
If a degree-$d$ refutation existed with $n\ge2^{2d}$,
Lemma~\ref{tx:lem:order-NS-restriction} would produce a chain refutation
forbidden by Proposition~\ref{tx:prop:order-induction-lower} at budget
$b=2d$. Therefore $n<2^{2d}$ for the finite minimum degree $d$,
which is the displayed inequality.
\end{proof}

\section{Why PC degree also tends to infinity}
We now use the fixed positive characteristic. The uniform comparison
theorem, Theorem~\ref{tx:thm:comparison}, supplies the
following input for the original order family: for each fixed PC
budget $d$ there is a finite NS budget $E(d)$, independent of $n$,
such that
\begin{equation}\label{tx:eq:order-comparison-input}
 d_{\PC}(\mathcal O_n)\le d
               \quad\Longrightarrow\quad
 d_{\NS}(\mathcal O_n)\le E(d).
\end{equation}
The function may grow quickly; only independence from $n$ is needed.

For a fixed $d$, inequality~\eqref{tx:eq:order-log-seed} eventually
exceeds $E(d)$. Consequently a degree-$d$ PC refutation is impossible
at every sufficiently large size. Since this holds for each fixed
$d$, the PC minimum degree tends to infinity too. No finite degree
can recur infinitely often in either proof system.

This uses comparison on $\mathcal O_n$, where the uniform hypothesis
holds. Applying it to the restricted chain would be invalid: that
nonuniform family has the degree-two PC proofs already displayed.
The comparison transfers \emph{qualitative unboundedness}; the linear
PC bound will come from lifting with its original PC degree unchanged.

\section{The same lifting mechanism gives the linear bound}
For this family, the invariant-target case of
Theorem~\ref{tx:thm:transfer} gives a particularly simple statement.
For either $S=\NS$ or $S=\PC$, if
\begin{equation}\label{tx:eq:order-lift-range}
 q=p^a>2d+3,\qquad n\ge\max\{9,4d+5\},
\end{equation}
then a degree-$d$ refutation at $n$ gives one at $n+q$. These numbers
come from $r=2d+3$ and the retained-label requirement $n\ge2r-1$.
We use the actual range of the transfer theorem; the profile chapter
also defines a slightly more conservative onset for general bookkeeping.

\begin{theorem}[A finite linear lower bound for the order family]
\label{tx:thm:order-linear}
Over every fixed field of prime characteristic, the NS and ordinary-PC
refutation degrees of $\mathcal O_n$ are finite at every positive
size, tend to infinity, and satisfy
\begin{equation}\label{tx:eq:order-linear-bound}
 d_{\PC}(\mathcal O_n)>\frac{n-5}{4},\qquad
 d_{\NS}(\mathcal O_n)>\frac{n-5}{4}\qquad(n\ge9).
\end{equation}
Both profiles are of type T2 in Chapter~\ref{tx:ch:profiles}.
\end{theorem}
\begin{proof}
Finiteness and positivity were proved in
Proposition~\ref{tx:prop:order-finite}. The preceding two sections show
that both minimum degrees tend to infinity.

Suppose a degree-$d$ refutation existed at a size satisfying
\eqref{tx:eq:order-lift-range}. Keep this $d$ fixed and choose one
admissible prime power $q$. Iterated forward lifting gives refutations
of degree at most $d$ at
\[
                            n,n+q,n+2q,\ldots.
\]
Every step remains eligible: the retained universe size grows, while $d,r,q$
stay fixed. This contradicts the fact that the minimum degree tends
to infinity. The contradiction works separately for NS and PC.

Now take any $n\ge9$ and let $d$ be its finite minimum degree in the
chosen system. If $n\ge4d+5$, the preceding contradiction applies.
Therefore $n<4d+5$, which rearranges to
$d>(n-5)/4$. Absence of any recurring finite degree is exactly the
T2 case of the later classification.
\end{proof}

The conclusion concerns unavoidable proof degree at every sufficiently
large size. It uses an infinite progression, not a drop in the degree:
the lifted proofs would all have the same budget $d$. A proof of
arbitrary length cannot evade the bound, because lifting transports
the complete degree-$d$ consequence space.

\section{What the later common projector must do for this example}
The application has now isolated the jobs of the algebraic construction.
The common projector will be a linear projection obtained from one
shared combination of label permutations. Write $\Lift$ for the
resulting linear map from smaller-universe polynomials to larger-universe ones.
The variable basis has two-label support; monomials of degree at most
$d$ use at most $2d$ labels. The usable-axiom source and the NS
multiplier--axiom source are covered by $r=2d+3$. Every closure stage
stays inside these fixed ambient spaces.

Consider the row input $\sum_{j\le n}F_{ij}-1$. Plain inclusion into
a larger universe leaves out its added summands. Setting all variables
involving deleted labels to zero instead turns a deleted row into
$-1$. The common projector must repair the actual axiom coefficient
map. Its axiom-lifting formula gives a linear combination of genuine
larger-instance inputs, rather than assuming that the same retained
polynomial remains an axiom.

For PC it must also preserve $V_j\cap P_{\le d-1}$ after cancellation,
and the joint tensor multiplication map. This ensures that every
lifted vector in the larger closure has a legal ordinary-PC derivation.
The target here is $1$, so the identity $\Lift(1)=1$ turns consequence
lifting into refutation lifting. The retained-label cost $2r-1=4d+5$
is what can be inverted to obtain the linear bound.

The separate comparison theorem uses a uniformly bounded chain of
invariant subspaces and an NS plateau to contain the whole PC closure.
In this example it rules out bounded PC degree before lifting supplies
the numerical lower bound. The restriction to the induction chain,
the clearance invariant, comparison, and the common projector each
have one task. Together they turn a sequential finite-order
contradiction into the same linear-degree mechanism used throughout
the book.
